\subsection*{\texttt{with-in} 块与 \texttt{where} 子句}

\co@with-in@ 将一个索引元组与形状范围绑定：

\begin{lstlisting}[language=croqtile,caption={索引 with-in 块}]
with index in [10, 10] {
  // index is ituple with 2 elements
}
\end{lstlisting}

\co@index@ 包含两个元素，也可以分别给元素命名：

\begin{lstlisting}[language=croqtile,caption={索引元组}]
with {x, y} in [10, 10] {...}
\end{lstlisting}

元组本身和各个元素可以同时命名：

\begin{lstlisting}[language=croqtile,caption={由元组导出的索引}]
with index = {x, y} in [10, 10] {...}
\end{lstlisting}

这里的 \co@index@ 是有界整数元组。

\co@with-in@ 只声明有界索引，不启动并行工作。需要遍历这些索引时，再使用 \co@foreach@。

\co@where@ 可以声明索引之间的对应关系：

\begin{lstlisting}[language=croqtile,caption={嵌套的有界组}]
with {m, n} in [M, N], {n_p, k} in [N_P, K] where n_p <-> n {
  // matmul implements with m,n,K. n_p is no long useful.
}
\end{lstlisting}

\co@n <-> n_p@ 约定两个索引在迭代中对应同一取值，因此可以在需要该索引的位置使用相应名字。
